\documentclass[10pt]{article}
\usepackage[T1]{fontenc}
\usepackage[utf8]{inputenc}
\usepackage[french]{babel}
\usepackage{newpxtext,newpxmath}
\usepackage[a4paper,top=1.5cm,bottom=1.35cm,left=1.5cm,right=1.5cm]{geometry}
\usepackage{microtype}
\usepackage{xcolor}
\usepackage{titlesec}
\usepackage[most]{tcolorbox}
\usepackage{enumitem}
\usepackage{array}
\usepackage{booktabs}
\usepackage{fancyhdr}

\definecolor{navy}{HTML}{1F3A5F}
\definecolor{grey}{HTML}{54606E}
\definecolor{green}{HTML}{2C6E49}
\definecolor{amber}{HTML}{9C6A16}
\definecolor{red}{HTML}{A6322E}
\definecolor{plum}{HTML}{7A4EA8}
\definecolor{boxgray}{HTML}{F3F5F8}
\definecolor{boxred}{HTML}{FCF2F1}
\definecolor{boxgreen}{HTML}{F1F7F3}
\definecolor{boxplum}{HTML}{F7F3FB}
\definecolor{boxcode}{HTML}{F7F8FA}
\definecolor{boxedge}{HTML}{D9E0EA}

\setlength{\parindent}{0pt}
\setlength{\parskip}{2.6pt}
\linespread{1.0}

\pagestyle{fancy}\fancyhf{}
\lhead{\small\color{grey}Scella}
\rhead{\small\color{grey}Feuille de route en modules --- 20 août 2026}
\cfoot{\small\color{grey}\thepage}
\renewcommand{\headrulewidth}{0.3pt}

\titleformat{\section}{\color{navy}\large\bfseries}{}{0pt}{}
\titlespacing{\section}{0pt}{11pt}{3pt}
\titleformat{\subsection}{\color{navy}\bfseries}{}{0pt}{}
\titlespacing{\subsection}{0pt}{7pt}{1.5pt}

\newtcolorbox{these}{colback=boxgray,colframe=boxedge,boxrule=0.5pt,arc=2.5pt,
  left=8pt,right=8pt,top=4.5pt,bottom=4.5pt}
\newtcolorbox{alerte}{colback=boxred,colframe=red!55,boxrule=0.8pt,arc=2.5pt,
  left=8pt,right=8pt,top=4.5pt,bottom=4.5pt}
\newtcolorbox{acquis}{colback=boxgreen,colframe=green!45,boxrule=0.8pt,arc=2.5pt,
  left=8pt,right=8pt,top=4.5pt,bottom=4.5pt}
\newtcolorbox{piste}{colback=boxplum,colframe=plum!45,boxrule=0.8pt,arc=2.5pt,
  left=8pt,right=8pt,top=4.5pt,bottom=4.5pt}
\newtcolorbox{extrait}{colback=boxcode,colframe=boxedge,boxrule=0.5pt,arc=1.5pt,
  left=7pt,right=7pt,top=3.5pt,bottom=3.5pt}

\newcommand{\best}[1]{\textcolor{green}{\textbf{#1}}}
\newcommand{\pire}[1]{\textcolor{red}{\textbf{#1}}}
\newcommand{\sz}[1]{\textcolor{grey}{\small\textsf{#1}}}
\newcommand{\code}[1]{\texttt{\small #1}}
\newcommand{\fl}{$\rightarrow$}
\newcommand{\mm}[1]{\textbf{\textcolor{navy}{#1}}}

\begin{document}

\begin{center}
{\Huge\bfseries\color{navy} Scella}\\[4pt]
{\large\color{grey}\itshape Finir C1 et C3 --- la feuille de route en modules}\\[4pt]
{\small\color{grey} 20 août 2026 \; · \; douze modules, chacun avec son critère d'acceptation vérifiable}
\end{center}
\vspace{2pt}

\begin{these}
\textbf{Pourquoi des modules plutôt qu'un calendrier.} Un calendrier suppose que rien ne bouge ; des
modules survivent au premier imprévu. Chacun de ceux qui suivent porte un \textbf{critère
d'acceptation vérifiable} --- une commande, pas une impression --- et ne dépend que de modules
nommés. Vous pouvez en faire deux et vous arrêter, ou changer l'ordre, sans réécrire le plan.

\textbf{Le chemin critique tient en quatre modules} : \mm{M0} le profil, \mm{M2} l'observation non
intrusive, \mm{M5} le pont vers le registre, \mm{M6} la conservation. Les huit autres sont utiles
mais ne bloquent personne.
\end{these}

\section{1 · Les douze modules}

\renewcommand{\arraystretch}{1.28}
\begin{center}\footnotesize
\begin{tabular}{@{}l >{\raggedright\arraybackslash}p{3.5cm} >{\raggedright\arraybackslash}p{6.6cm} >{\raggedright\arraybackslash}p{2.0cm} c@{}}
\toprule
& \textbf{Module} & \textbf{Fini quand} \sz{(vérifiable)} & \textbf{Dépend de} & \sz{Taille} \\
\midrule
\multicolumn{5}{@{}l@{}}{\textit{\textcolor{grey}{— Socle —}}} \\
\mm{M0} & \textbf{Le profil servi} & un clone frais $+$ le profil documenté sert un reçu \textbf{non vide} & --- & \sz{3 j} \\
\mm{M1} & Le garde unique & \code{grep -rln 'fn is\_local\_endpoint'} renvoie \textbf{un} fichier & --- & \sz{2 j} \\
\addlinespace[3pt]
\multicolumn{5}{@{}l@{}}{\textit{\textcolor{green}{— Voie C3 : le coin —}}} \\
\mm{M2} & \textbf{L'observation non intrusive} & on \textbf{tue} la passerelle en plein run : le client reçoit quand même sa réponse & M0 & \sz{1--2 sem} \\
\mm{M3} & La panne qui ne coupe rien & un test tue le SLM et le store : l'appel passe, \code{degrade\_reason} estampillé & M2 & \sz{3 j} \\
\mm{M4} & L'installation en une commande & machine vierge, \code{docker compose up} \fl{} un reçu re-prouvé, \textbf{sans clé} & M0, M2 & \sz{1 sem} \\
\addlinespace[3pt]
\multicolumn{5}{@{}l@{}}{\textit{\textcolor{plum}{— La jonction —}}} \\
\mm{M5} & \textbf{Le pont vers le registre} & un run de la passerelle produit un dossier \code{reg-evidence} \textbf{sans geste manuel} & M2 & \sz{1--2 sem} \\
\addlinespace[3pt]
\multicolumn{5}{@{}l@{}}{\textit{\textcolor{navy}{— Voie C1 : l'expansion —}}} \\
\mm{M6} & \textbf{La conservation limitée} & un cas expiré disparaît, et \textbf{tout le reste se re-prouve encore} & M5 & \sz{1 sem} \\
\mm{M7} & Le service de lecture & \code{audit-api} répond en HTTP sur un port \sz{(aujourd'hui : bibliothèque)} & M0 & \sz{1 sem} \\
\mm{M8} & La console opposable & auth $+$ TLS $+$ un tenant ne voit \textbf{jamais} l'autre & M7 & \sz{1 sem} \\
\addlinespace[3pt]
\multicolumn{5}{@{}l@{}}{\textit{\textcolor{amber}{— L'opposabilité —}}} \\
\mm{M9} & L'ancre tierce (3b) & \code{AnchorProof} de méthode \code{rfc3161} produite et vérifiée & M4 & \sz{2 sem} \\
\addlinespace[3pt]
\multicolumn{5}{@{}l@{}}{\textit{\textcolor{grey}{— Hors chemin —}}} \\
\mm{MZ} & \textbf{La mesure du Cadrage} & le plafond du compresseur et le taux de préservation sont \textbf{chiffrés} & --- & \sz{1,5 j} \\
\mm{MX} & Le collecteur ouvert & un client Python écrit un nœud scellé sans toucher à Rust & M7 & \sz{3--4 sem} \\
\bottomrule
\end{tabular}
\end{center}

\sz{\mm{MZ} est indépendant, local et \textbf{gratuit} : faisable n'importe quand, il peut tuer le
Cadrage sans rien coûter. \mm{MX} est \textbf{conditionnel} — ne le faites que si trois prospects
demandent le registre sans vouloir router leur trafic.}

\section{2 · Les quatre modules qui méritent un mot}

\subsection{2.1 · \mm{M2} --- l'observation non intrusive \sz{(le déblocage)}}

C'est le seul module dont dépend la vente. Aujourd'hui la passerelle est dans le \textbf{chemin
critique} : si elle tombe, l'accès aux modèles tombe. Aucune équipe plateforme d'assureur n'accepte
cela d'un fournisseur inconnu.

En mode observation, la passerelle \emph{voit} le trafic et le scelle, mais ne le \emph{porte} pas.
Le risque d'adoption tombe alors au niveau de C1 --- c'est-à-dire à zéro.

\begin{acquis}
\textbf{Le critère d'acceptation est brutal, et c'est voulu :} on tue le processus de la passerelle
au milieu d'un run, et le client reçoit quand même sa réponse. Tant que ce test n'existe pas, le
module n'est pas fini --- et la promesse « sans risque » est une affirmation, pas une propriété.
\end{acquis}

\subsection{2.2 · \mm{M5} --- le pont vers le registre \sz{(la valeur)}}

C'est le module qui vaut le plus cher du plan, et il est court. Il transforme C1 d'une vente froide
de dix-huit mois en un \textbf{upsell sur un client déjà déployé}.

Le substrat existe : \code{provenance-bridge} traduit déjà un run en nœuds de la pile B, et la route
\code{bundle} les sert. Ce qui manque est le dernier maillon --- que ces nœuds nourrissent les
gabarits de \code{reg-evidence} sans intervention.

\subsection{2.3 · \mm{M6} --- la conservation limitée \sz{(le risque levé)}}

Je le croyais difficile : purger dans une chaîne scellée casse la re-preuve. \textbf{Vérification
faite, ce n'est pas le cas ici.}

\begin{acquis}
Les taps ne stockent que des \textbf{empreintes} --- \code{inputs\_hash}, \code{query\_hash},
\code{content\_hash} --- et un test nommé \code{hash\_is\_deterministic\_and\_minimizes} le fige.
\textbf{L'effacement RGPD et la chaîne inforgeable ne s'opposent donc pas par construction} : le
contenu vit chez le client, la chaîne ne porte que des empreintes.
\end{acquis}

Reste une règle de conception à respecter, et elle est simple : \textbf{purger par \emph{cas} entier,
jamais par nœud.} Chaque cas a sa propre chaîne et sa propre racine ; supprimer un cas complet laisse
tous les autres re-prouvables. Retirer un nœud au milieu d'une chaîne la casserait --- c'est
exactement l'attaque par troncature, faite par soi-même.

\begin{alerte}
\textbf{La tension à arbitrer n'est pas technique, elle est juridique.} L'AI Act art.\,12 impose de
\emph{conserver} les journaux ; le RGPD art.\,5 impose de \emph{limiter} la conservation. La politique
est donc un \textbf{intervalle}, pas une durée --- et c'est un juriste qui le fixe, pas vous. Faites
relire avant de coder la purge.
\end{alerte}

\subsection{2.4 · \mm{M9} --- l'ancre tierce \sz{(un raccordement, pas une invention)}}

\code{forensic-export} porte déjà \code{AnchorMethod::Rfc3161}, une \code{AnchorProof} structurée et
un \code{verify\_anchor} qui exécute le contrôle décisif --- la liaison de racine. Ce qui manque est
d'en \textbf{produire} une : aujourd'hui seule la méthode \code{Declared} sort du pont.

C'est ce module qui fait passer le reçu de \emph{tamper-evident pour un détenteur de clé} à
\textbf{opposable au détenteur lui-même}.

\section{3 · L'ordre, et les deux jalons}

\begin{extrait}
\small
\code{M0 -> M2 -> M3} \hfill \sz{la panne qui ne coupe rien} \\
\code{M0 -> M2 -> M4 -> M9} \hfill \sz{l'installation, puis l'opposabilité} \\
\code{M0 -> M2 -> M5 -> M6} \hfill \sz{\textbf{le chemin critique}} \\
\code{M0 -> M7 -> M8} \hfill \sz{la surface de lecture} \\[3pt]
\code{M1} \sz{et} \code{MZ} \sz{: indépendants, faisables n'importe quand \quad ·\quad}
\code{MX} \sz{: conditionnel}
\end{extrait}

\begin{center}\small
\renewcommand{\arraystretch}{1.25}
\begin{tabular}{@{}>{\raggedright\arraybackslash}p{4.6cm} >{\raggedright\arraybackslash}p{5.0cm} >{\raggedright\arraybackslash}p{6.0cm}@{}}
\toprule
\textbf{Jalon} & \textbf{Atteint après} & \textbf{Ce qu'il autorise} \\
\midrule
\best{C3 vendable} & \mm{M0} $+$ \mm{M2} $+$ \mm{M3} $+$ \mm{M4} & chercher le premier client \\
\textcolor{navy}{\textbf{C1 livrable}} & $+$ \mm{M5} $+$ \mm{M6} $+$ \mm{M7} $+$ \mm{M8} & l'upsell sur ce client \\
\bottomrule
\end{tabular}
\end{center}

\begin{alerte}
\textbf{Le premier jalon est le seul qui décide.} S'il est atteint et qu'aucun prospect n'installe la
passerelle, le problème n'est \textbf{pas} le produit --- c'est la distribution, et coder \mm{M5} à
\mm{M9} n'y changera rien. Arrêtez-vous là et cherchez pourquoi.
\end{alerte}

\section{4 · Ce qui reste exclu}

\begin{itemize}[leftmargin=1.2em,itemsep=2pt,topsep=2pt]
\item \textbf{Le runtime nano} --- six crates, aucun appelant. Le retirer, pas le porter.
\item \textbf{La lignée latente comme produit} --- gardez le canari de Lipschitz comme lettre de
      créance, pas comme catalogue.
\item \textbf{Fil C et l'essaim médiateur} --- feuille de route, pas produit.
\item \textbf{Le Cadrage} --- jusqu'à \mm{MZ}. Il est gratuit : faites-le avant d'en parler à quiconque.
\item \textbf{\mm{MX}} --- jusqu'à ce que trois prospects le demandent \emph{et} le paient.
\end{itemize}

\begin{these}
\small\textbf{Rappel de cadrage, inchangé dans toute la série.} Le reçu atteste l'\textbf{intégrité}
d'une suite de pas, jamais la \textbf{correction} d'une décision. Il est \emph{tamper-evident pour un
détenteur de clé}, pas \emph{trustless pour un tiers} --- et c'est \mm{M9} qui corrige cela. Les
participants du mode réel exécutent une fusion latente Mode-1 numérique : ce ne sont pas des « vrais
SLM ». Les \textbf{trous} et les \textbf{critères} de ce document sont vérifiables ; les
\textbf{tailles} sont des estimations.
\end{these}

\section{Annexe · Vérifier l'état de départ}

\renewcommand{\arraystretch}{1.22}
\begin{center}\small
\begin{tabular}{@{}>{\raggedright\arraybackslash}p{5.0cm} >{\raggedright\arraybackslash}p{11.0cm}@{}}
\toprule
\textbf{Le trou que le module comble} & \textbf{Ce qui le montre aujourd'hui} \\
\midrule
\mm{M0} aucun profil & \code{grep -rl EFE\_META\_PROVENANCE . -{}-include='*.md' -{}-include='*.yml'} \sz{\fl{} vide} \\
\mm{M1} garde dupliqué & \code{grep -rln 'fn is\_local\_endpoint' crates -{}-include='*.rs'} \sz{\fl{} 2 fichiers} \\
\mm{M2} pas d'observation & \code{grep -rniE 'observe\_only|shadow|fail\_open' crates/efe-swarm-gateway/src} \sz{\fl{} 0} \\
\mm{M6} pas de rétention & \code{grep -rniE 'retention|purge|ttl' crates/audit-api/src crates/provenance-capture/src} \sz{\fl{} 0} \\
\mm{M6} la chaîne ne porte que des empreintes & \code{grep -n 'hash\_hex' crates/provenance-capture/src/tap.rs} \\
\mm{M7} \code{audit-api} est une bibliothèque & \code{ls crates/audit-api/src/bin} \sz{\fl{} absent} \\
\mm{M9} seul \code{Declared} est produit & \code{grep -rn 'Rfc3161' crates/provenance-bridge/src} \sz{\fl{} vide} \\
\bottomrule
\end{tabular}
\end{center}

Un avertissement pour finir. Les \textbf{critères d'acceptation} de ce document sont écrits pour être
exécutables : si l'un d'eux ne peut pas l'être, c'est le critère qu'il faut réécrire, pas le module
qu'il faut déclarer fini. Et le seul jalon qui ne dépende pas de votre vitesse est le premier ---
parce qu'il attend la réponse de quelqu'un d'autre.

\end{document}
